% DRAFT (Claude, 2026-11-10, Anders): 5.2 written with the literature-review skill.
% Plan and exemplar observations in PLAN-5-2.md; every citation's passage in CHANGES-5-2.md.
% Fourth version after three independent review rounds.
% Plan: (1) CLAIM two perspectives, system models optimise the whole system and produce prices
%       (2) CLAIM the system's problem is not the producer's (Helseth 2023)
%       (3) GROUP owner studies: the gain from pumping depends on the case
%       (4) COMPARE price-maker representations and their limits
%       (5) CLAIM why price-taking is assumed
%       (6) CRITIQUE its limit for a large plant
%       (7) SYNTH, ending on what is left open for a study that varies capacities

\subsection{Valuation Perspective}
\label{sec:lit_perspective}

Hydropower scheduling is studied from two perspectives, and the perspective is
crucial for what a result can say about an investment. System studies model
several areas with generation and demand, connected by transmission, and
optimise the system as a whole under one objective. \citet{helseth2013}
schedule the Icelandic hydro-thermal system in this way.
\citet{henden2016economic} and \citet{helseth2017} value pumped-storage
expansions in the Nordic system by the change in socio-economic surplus.
% src: helseth2013 pp. 1427, 1431; henden2016economic pp. 117-118; helseth2017 pp. 441, 445
Prices are a result of the simulation in such models
\citep{henden2016economic}.
% src: henden2016economic p. 118 ("Simulation results are power prices, reservoir level, exchanges")
The result is therefore a system-wide dispatch and its prices, not a strategy
for one plant.

However, the system's problem is not the producer's. In the deregulated
Nordic market a producer has no obligation to cover demand, and
\citet{helseth2023hydropower} describe how the scheduling problem is instead
distributed to each producer.
% src: helseth2023hydropower p. 7
Each producer maximises its own expected profit given a forecast of future
prices, and the system models supply those forecasts rather than the
decisions.
% src: helseth2023hydropower pp. 7-8
Moreover, the surplus a system model maximises is shared between consumers,
producers and the grid, so \citet{henden2016economic} compute the owners'
share separately.
% src: henden2016economic pp. 119-120, eqs. 1-5
A favourable system-level result for pumped storage may thus say little about
what its owner would earn.

Owner-oriented studies place the revenue or profit of a particular asset in
the objective. The gains from pumping they report vary with the price year
and between cases. \citet{ak2019quantifying} operate an existing Turkish
cascade as pumped storage and report a net revenue gain whose size varies
with the price year.
% src: ak2019quantifying pp. 739-740 (existing cascade), p. 747, Table 2
A pump added to a Norwegian cascade with environmental restrictions gave only
a small gain in its case \citep{alic2023optimal}.
% src: alic2023optimal pp. 16-22 (constraints), p. 30 ("only" +2.35 %)
Together these studies suggest that pumping can raise an asset's profit, but
that the gain found for one plant may not carry over to another.

A further distinction within the owner perspective is whether the producer
influences prices. This depends on the size of its bids relative to the
market, so the question arises particularly in the small reserve markets
\citep{perezdiaz2015trends,perezdiaz2020medium}.
% src: perezdiaz2020medium p. 3; perezdiaz2015trends p. 772
\citet{abgottspon2013stochastic} let prices depend linearly on the
quantities their pumped-storage producer trades, with a relationship
estimated from experience.
% src: abgottspon2013stochastic pp. 1, 4
Instead, \citet{perezdiaz2020medium} model price-making in the reserve
market only, which in their case changes expected profit very little.
% src: perezdiaz2020medium pp. 7-8
However, such a formulation rests on an assumed market response and makes the
revenue function non-convex, which the long-horizon decomposition methods
handle only with extensions \citep{fullnerRebennack2025}.
% src: fullnerRebennack2025 p. 460

Many owner-oriented studies assume a price-taker, for whom prices are
uncertain inputs unaffected by its own decisions
\citep{helseth2016,ak2019quantifying,alic2023optimal,bore2026assessing}.
% src: helseth2016 p. 934; ak2019quantifying p. 741; alic2023optimal p. 24; bore2026assessing p. 3
In these studies the assumption is justified in three ways.
\citet{helseth2023hydropower} regard it as a condition for an efficient
market that holds fairly well in the Nordic market.
% src: helseth2023hydropower p. 7
Further, \citet{bore2026assessing} and \citet{ak2019quantifying} argue from
the size of the plant relative to its price zone. \citet{sioshansi2009estimating}
add the practical reason that it keeps the scheduling problem linear.
% src: bore2026assessing p. 3; ak2019quantifying p. 741; sioshansi2009estimating p. 270
The size-based reasons hold only for a plant that is small relative to its
market.

The limitation of price-taking is that a large plant moves the prices it
trades on. \citet{dallinger2019socioeconomic} find in their system model
that the price spread available to pumped storage narrows as storage
capacity grows. Similarly, \citet{henden2016economic} find that a large
expansion in Norway largely removes its own profitability.
% src: dallinger2019socioeconomic pp. 488-489, Table 1; henden2016economic p. 122
Consequently, a price-taking model may overstate the value of a large
expansion, which is especially relevant when a plant's capacities are varied.

Overall, the operation of a particular asset is studied from the owner's
perspective, with the expected profit of that asset as the objective and
prices as uncertain inputs. A price-taking formulation is a common choice
when the plant is small relative to its markets, and system studies are read
alongside it for what they reveal about prices. What the reviewed studies
leave open is how far the price effect of a large plant changes the value of
its capacities, which is particularly relevant for a study that varies them.
